% Findings and next steps

\section{Conclusion}
\label{sec:conclusion}

The architecture developed here puts DataOps and MLOps on one Kubernetes control plane using open source tools throughout. The nine layers derive from convergence across sources with different methods, components come from scored comparison, and the disciplines join at the three points of Section~\ref{subsec:bridges}. The value is that integration effort previously repeated per project is consolidated into a single declarative description, so rebuilding no longer depends on one engineer's undocumented knowledge.

The quality gate rejecting bad batches and the automatic firing of the retraining chain are the two results that met an acceptance criterion. The rebuild of the whole platform from one Git source is the strongest of the structural results, and most requirements sit at that tier, which establishes that the control plane forms as designed, not that it performs at production load.

The most immediate next step is load characterisation under real traffic, because behaviour under load is unknown. Drift observation needs longer runs on real streams before detection accuracy can be assessed, and the per-entity comparison of online against offline feature values should be automated so that the equality argued from construction is confirmed by measurement.
